% MA 214 Discussion 11 -- covers Lecture 16, and doubles as Quiz 3 / Module 3 review.
%
% Student handout. Page 1 is the Module 3 concept review; the question set runs from page 2.
% Questions are tagged with the Module 3 objective they test.
%
% Compile from inside this folder:  pdflatex Discussion11.tex

\documentclass[11pt]{article}
\usepackage[margin=1in]{geometry}
\usepackage{parskip}
\usepackage{fancyhdr}
\usepackage{array}
\usepackage{booktabs}
\usepackage{enumitem}
\usepackage{amsmath}
\usepackage{tabularx}
\usepackage{listings}

\pagestyle{fancy}
\fancyhf{}
\cfoot{\thepage}
\rhead{Discussion 11}
\lhead{MA 214: Applied Statistics}
\renewcommand{\headrulewidth}{.4mm}

\newcommand{\blankline}[1]{\underline{\hspace{#1}}}
\newcolumntype{L}{>{\raggedright\arraybackslash}X}

\lstset{
  basicstyle=\ttfamily\small,
  columns=fullflexible,
  frame=single,
  framerule=0.4pt,
  breaklines=true,
  showstringspaces=false,
  aboveskip=8pt,
  belowskip=8pt
}

% Section header: bold title over a hairline rule.
\newcommand{\parttitle}[1]{%
  \par\medskip\noindent
  {\large\bfseries #1}\par
  \vspace{-.5em}\noindent\rule{\linewidth}{0.4pt}\par\smallskip}

% Writing space.
\newcommand{\answerspace}[1]{\vspace{#1}}

% Objective tag printed after a question title.
\newcommand{\lo}[1]{{\normalfont\small (#1)}}

% Flush-left question list: the label runs in at the margin and the body wraps
% back to the margin, so nothing is pushed far to the right.
\newlist{questions}{enumerate}{1}
\setlist[questions]{label=\textbf{Question \arabic*.}, wide=0pt, itemsep=1.4em, topsep=.8em}

\newlist{parts}{enumerate}{1}
\setlist[parts]{label=\textbf{(\Alph*)}, leftmargin=1.6em, labelsep=.45em,
  itemsep=.7em, topsep=.5em}

\begin{document}

\noindent\textbf{Name:}\blankline{2.3in}\hfill\textbf{BUID:}\blankline{1.6in}

\begin{center}
  {\Large\bfseries Discussion 11: Logistic Inference, and All of Module 3}
\end{center}

\noindent\textbf{Goals:} run Module 3's last tool and then the whole module. Each question is
labelled with the objective it tests, so you can see which one to go back to before the quiz.

\parttitle{Part 1: Module 3 Concept Review}

{\small

\noindent The whole module on one page. Everything in Part 2 can be answered from it.

\begin{enumerate}[label=\textbf{\arabic*.}, wide=0pt, itemsep=.4em, topsep=.5em]

\item \textbf{M3.L01 Inference for a slope.} $\hat\beta_1$ is a statistic. \emph{Road A:} shuffle the
pairing of $y$ to $x$ to test $\beta_1 = 0$; bootstrap whole $(x,y)$ pairs for an interval.
\emph{Road B:} $t = \hat\beta_1/\text{SE}$ with df $= n-p-1$, and
$\hat\beta_1 \pm t^{\ast}\text{SE}$. To test against a value other than zero, use
$t = (\hat\beta_1 - c)/\text{SE}$, or ask whether $c$ lies in the interval.

\item \textbf{M3.L02 Conditions and remedies.} Road B needs \textbf{LINE}: linear mean, independent
observations, normal residuals, equal variance. Check L and E on residuals against fitted values, N
on a Q--Q plot. A funnel means E failed, a curve means L failed. Remedy by transforming, then
\emph{re-check}. Agreement between the two roads is not evidence that the conditions hold.

\item \textbf{M3.L03 Reading multiple regression output.} Every coefficient is a change in $y$ per
unit of that predictor \emph{holding the others fixed}. Collinear predictors inflate standard errors,
so each can look useless while the pair matters; interpretation breaks, prediction survives.

\item \textbf{M3.L04 Interactions and transformations.} With an indicator $g$, the slope in $x$ is
$\hat\beta_1$ when $g=0$ and $\hat\beta_1+\hat\beta_3$ when $g=1$, so the interaction coefficient is
a \emph{difference of slopes}. Once an interaction is present, $g$'s main effect is only the gap at
$x=0$. A log--log slope is an elasticity.

\item \textbf{M3.L05 Prediction and cross-validation.} At $x^{\ast}$,
\[
\text{CI (mean): } \hat{y} \pm t^{\ast}\text{SE}_{\text{fit}},
\qquad
\text{PI (one case): } \hat{y} \pm t^{\ast}\sqrt{\hat\sigma^2 + \text{SE}_{\text{fit}}^2}.
\]
More data shrinks the CI toward nothing but never takes the PI below about
$\pm t^{\ast}\hat\sigma$. $R^2$ always prefers the biggest model; adjusted $R^2$ and AIC charge for
parameters but are still computed on the training data. Only cross-validation predicts each point
from a model that never saw it.

\item \textbf{M3.L06 Logistic inference.} Coefficients are on the log-odds scale, so
$\text{OR} = e^{\hat\beta}$ and $e^{c\hat\beta}$ for a $c$-unit change. Build every interval on the
\textbf{log-odds scale first}, then transform: exponentiate the endpoints for an odds ratio, or apply
$p = e^{\eta}/(1+e^{\eta})$ for a probability. The result is asymmetric, and that is correct;
building a symmetric interval directly on the probability scale can produce impossible values below
$0$ or above $1$. Judge predictions with accuracy \emph{next to the base rate}, plus the Brier score
$\overline{(\hat{p}-y)^2}$ and the log-loss
$-\overline{y\log\hat{p} + (1-y)\log(1-\hat{p})}$; log-loss punishes confident and wrong far harder.

\item \textbf{Four traps.} Averaging or halving probabilities where you should work in log odds.
Reading a main effect's $p$-value as ``the group does not matter'' in an interaction model. Reporting
a CI when a single new case was asked about. Quoting in-sample accuracy as if it were out-of-sample.

\end{enumerate}

}

\newpage

\parttitle{Part 2: Question Set}

\noindent\textbf{Scenario.} For $n = 120$ branches, whether the branch \textbf{met} its annual
target is modelled on annual \texttt{visits} (thousands) and \texttt{type} (Branch, Central, Kiosk;
Branch is the reference). $67$ branches missed the target and $53$ met it.

\begin{center}
\renewcommand{\arraystretch}{1.2}
\begin{tabular}{lccccc}
\toprule
 & \textbf{Estimate} & \textbf{SE} & \textbf{$p$} & \textbf{OR} & \textbf{95\% CI for OR} \\
\midrule
(Intercept) & $-6.1881$ & 1.0969 & $<0.001$ & --- & --- \\
\texttt{visits} & $0.0739$ & 0.0132 & $<0.001$ & 1.077 & $[1.049,\; 1.105]$ \\
\texttt{type}Central & $1.4141$ & 0.6359 & 0.026 & 4.113 & $[1.183,\; 14.303]$ \\
\texttt{type}Kiosk & $-0.0798$ & 0.6029 & 0.895 & 0.923 & $[0.283,\; 3.010]$ \\
\bottomrule
\end{tabular}
\end{center}

\begin{questions}

%%%%%%%%%%%%%%%%%%%%%%%%%%%%%%%%%%%%%%
\item \textbf{Reading the table.} \lo{M3.L03, M3.L06}

\begin{parts}
\item The \texttt{visits} odds ratio is $1.077$ per thousand. Give the odds ratio for ten thousand
more visits, using $e^{0.739} = 2.093$, and interpret it in one sentence.
\answerspace{4em}

\item The interval for the \texttt{type}Central odds ratio is $[1.183,\; 14.303]$, which is wildly
lopsided about $4.113$. Explain why an odds-ratio interval comes out asymmetric even though the
interval it came from was symmetric.
\answerspace{4em}

\item The \texttt{type}Kiosk interval is $[0.283,\; 3.010]$. What does it mean that this interval
contains $1$, and what would you tell the director about kiosks?
\answerspace{4em}
\end{parts}

%%%%%%%%%%%%%%%%%%%%%%%%%%%%%%%%%%%%%%
\item \textbf{An interval for a predicted probability.} \lo{M3.L06} For a \textbf{Central} branch
with \texttt{visits} $= 70$, the model gives $\hat\eta = 0.3975$ with
$\text{SE}(\hat\eta) = 0.4628$. You may use $e^{0.3975} = 1.488$, \; $e^{-0.5095} = 0.601$, \;
$e^{1.3045} = 3.686$.

\begin{parts}
\item Compute $\hat{p}$ for this branch.
\answerspace{3em}

\item Build the 95\% interval on the \textbf{log-odds} scale, then convert both endpoints to
probabilities.
\answerspace{5em}

\item How far does your interval reach below $\hat{p}$, and how far above? Why is the asymmetry
correct rather than a mistake?
\answerspace{4em}

\item A classmate skips the log-odds step and builds a symmetric interval directly on the
probability scale. For a \textbf{Kiosk} branch with $\texttt{visits} = 20$, where
$\hat{p} = 0.0082$, their method returns $[-0.0066,\; 0.0231]$ while the correct method returns
$[0.0014,\; 0.0485]$. Name the fatal defect in their answer, and state the rule it violates.
\answerspace{4.5em}
\end{parts}

%%%%%%%%%%%%%%%%%%%%%%%%%%%%%%%%%%%%%%
\item \textbf{Judging the predictions.} \lo{M3.L05, M3.L06} At a threshold of $0.5$ the model
predicts correctly for $97$ of the $120$ branches, with sensitivity $0.792$ and specificity
$0.821$.

\begin{parts}
\item Compute the accuracy, then the base rate for this data set. Does the model beat it?
\answerspace{4em}

\item Complete the loss table for four illustrative branches.

\begin{center}
\renewcommand{\arraystretch}{1.6}
\begin{tabular}{ccccc}
\toprule
$\hat{p}$ & $y$ & $(\hat{p}-y)^2$ & $-\log$ of the probability given to the truth \\
\midrule
0.90 & 1 & \blankline{.9in} & \blankline{1.1in} \\
0.70 & 1 & \blankline{.9in} & \blankline{1.1in} \\
0.30 & 0 & \blankline{.9in} & \blankline{1.1in} \\
0.20 & 1 & \blankline{.9in} & \blankline{1.1in} \\
\bottomrule
\end{tabular}
\end{center}

\noindent You may use $\log 0.9 = -0.105$, $\log 0.7 = -0.357$, $\log 0.2 = -1.609$.

\item The last branch was given $0.20$ and met its target. Compare how hard the two losses punish
that one case, and say which loss you would use if confident wrong answers are especially costly.
\answerspace{4em}

\item In-sample the Brier score is $0.131$ and the log-loss $0.415$; under 10-fold
cross-validation they are $0.147$ and $0.471$, and accuracy falls from $0.808$ to $0.775$. Why do
all three get worse, and which set of numbers do you report?
\answerspace{4.5em}
\end{parts}

%%%%%%%%%%%%%%%%%%%%%%%%%%%%%%%%%%%%%%
\item \textbf{The linear side of Module 3.} \lo{M3.L01, M3.L02, M3.L04} Short answers.

\begin{parts}
\item A residual-versus-fitted plot for a linear model fans out to the right. Which LINE condition
failed, what do you do about it, and what must you do after that?
\answerspace{4em}

\item A model has $\hat\beta_{\texttt{hours}} = 1.03$ and
$\hat\beta_{\texttt{hours}:\texttt{Central}} = 0.96$. Give the slope for each group, and say what
the second coefficient is a difference of.
\answerspace{4em}

\item You want a plausible range for the slope but the residuals are clearly not normal. Which road
do you take, and what exactly do you resample?
\answerspace{3.5em}

\item A slope is $1.35$ with SE $0.115$ and a 95\% interval of $[1.119,\; 1.591]$. Without further
arithmetic, what do you conclude about $H_0: \beta_1 = 1.5$?
\answerspace{3.5em}
\end{parts}

%%%%%%%%%%%%%%%%%%%%%%%%%%%%%%%%%%%%%%
\item \textbf{Side by side.} \lo{all of Module 3} Complete the comparison.

\begin{center}
\renewcommand{\arraystretch}{2}
\begin{tabularx}{\linewidth}{LLL}
\toprule
 & \textbf{Linear regression} & \textbf{Logistic regression} \\
\midrule
The response is & & \\
The model is linear in & & \\
Fitted by & & \\
One coefficient means & & \\
Interval for a coefficient & & \\
Conditions are checked by & & \\
Predictions are judged by & & \\
\bottomrule
\end{tabularx}
\end{center}

\begin{parts}
\item One row differs more than the others between the two columns. Which, and why is that the
row where students most often go wrong?
\answerspace{3em}

\item A classmate's write-up ends with the two sentences below. Mark each \textbf{defensible} or
\textbf{not}, and rewrite the ones that are not.
\begin{enumerate}[label=\textbf{S\arabic*.}, leftmargin=2.2em, itemsep=.25em, topsep=.35em]
\item ``Central branches are $4.11$ times as likely to meet their target as ordinary branches.''
\item ``Cross-validation gave $77.5\%$ accuracy against $80.8\%$ in-sample, so the
cross-validated figure is too pessimistic and we report the in-sample one.''
\end{enumerate}
\answerspace{5em}
\end{parts}

\end{questions}

\end{document}
